% META: {"id": "TNSI-PRATIQUE-P8", "chapitre": "TNSI-STRUCTURES-DONNEES", "type_objet": "banque_pratique", "capacites": ["T-STRUCT-05C", "T-STRUCT-05D"], "capacites_codes": ["C13", "C14"], "duree_min": 60, "autorite": "MENE2516123N", "mode_creation": "ex_nihilo", "sources_inspiration": [], "notice": "ORIGINAL_NEXUS_TRAINING_MATERIAL_NOT_OFFICIAL_EXAM_CONTENT", "status": "generated"}
% BEGIN-VERIFY
% def matrice_vers_listes(matrice):
%     """Precondition : matrice carree de 0 et de 1.
%     Postcondition : dict sommet -> liste croissante de successeurs."""
%     n = len(matrice)
%     assert all(len(ligne) == n for ligne in matrice), "matrice non carree"
%     assert all(v in (0, 1) for ligne in matrice for v in ligne), "valeurs hors 0/1"
%     return {i: [j for j, v in enumerate(ligne) if v == 1]
%             for i, ligne in enumerate(matrice)}
% def listes_vers_matrice(listes):
%     """Precondition : les successeurs sont des sommets declares."""
%     n = len(listes)
%     assert all(0 <= j < n for s in listes.values() for j in s), "successeur inconnu"
%     matrice = [[0] * n for _ in range(n)]
%     for i, successeurs in listes.items():
%         for j in successeurs:
%             matrice[i][j] = 1
%     return matrice
% # jeu de tests remis avec le sujet
% m = [
%     [0, 1, 1, 0],
%     [0, 0, 0, 1],
%     [0, 0, 0, 1],
%     [0, 0, 0, 0],
% ]
% listes = matrice_vers_listes(m)
% assert listes == {0: [1, 2], 1: [3], 2: [3], 3: []}
% assert listes_vers_matrice(listes) == m
% # aller-retour sur plusieurs matrices, dont les cas extremes
% vide = [[0] * 4 for _ in range(4)]
% complet = [[0 if i == j else 1 for j in range(4)] for i in range(4)]
% boucle = [[1, 0], [0, 1]]
% for essai in (m, vide, complet, boucle):
%     assert listes_vers_matrice(matrice_vers_listes(essai)) == essai
% # les sommets isoles gardent une liste vide, et ne disparaissent pas
% assert matrice_vers_listes(vide) == {0: [], 1: [], 2: [], 3: []}
% assert len(matrice_vers_listes(vide)) == 4
% # comptages
% def aretes_matrice(matrice):
%     return sum(sum(ligne) for ligne in matrice)
% def aretes_listes(listes):
%     return sum(len(s) for s in listes.values())
% for essai in (m, vide, complet, boucle):
%     assert aretes_matrice(essai) == aretes_listes(matrice_vers_listes(essai))
% assert aretes_matrice(m) == 4
% assert aretes_matrice(complet) == 12
% assert aretes_matrice(boucle) == 2      # deux boucles sur soi-meme
% # les preconditions sont refusees
% for mauvais in (lambda: matrice_vers_listes([[0, 1], [0]]),
%                 lambda: matrice_vers_listes([[0, 2], [0, 0]]),
%                 lambda: listes_vers_matrice({0: [5], 1: []})):
%     refuse = False
%     try:
%         mauvais()
%     except AssertionError:
%         refuse = True
%     assert refuse
% END-VERIFY
\section*{Situation pratique P8 --- passer d'une matrice d'adjacence aux listes de successeurs}

\textit{Durée : 1 heure. Travail sur machine, à partir du document fourni.}

\medskip
\textit{Contenu original Nexus --- ce n'est pas un sujet officiel d'examen.}

\subsection*{Document fourni}

Le fichier \lstinline{graphes.py} contient les deux spécifications et le jeu de
tests.

\begin{python}
def matrice_vers_listes(matrice):
    """Precondition : matrice carree de 0 et de 1.
    Postcondition : dict sommet -> liste croissante de successeurs."""
    ...

def listes_vers_matrice(listes):
    """Precondition : les successeurs sont des sommets declares."""
    ...
\end{python}

La matrice de travail :

\begin{center}
\begin{tabular}{|c|c|c|c|c|}
\hline
 & $0$ & $1$ & $2$ & $3$ \\
\hline
$0$ & $0$ & $1$ & $1$ & $0$ \\
\hline
$1$ & $0$ & $0$ & $0$ & $1$ \\
\hline
$2$ & $0$ & $0$ & $0$ & $1$ \\
\hline
$3$ & $0$ & $0$ & $0$ & $0$ \\
\hline
\end{tabular}
\end{center}

\subsection*{Travail demandé}

\begin{enumerate}
  \item Compléter les deux fonctions. Un sommet sans successeur doit apparaître
    dans les listes, avec une liste \textbf{vide}.

  \item Rendre les préconditions exécutables : matrice non carrée, valeur autre
    que $0$ ou $1$, successeur hors des sommets déclarés. Vérifier que les trois
    sont refusées.

  \item Vérifier l'\textbf{aller-retour} sur quatre graphes : celui de l'énoncé,
    un graphe sans aucune flèche, un graphe complet à quatre sommets, et un
    graphe à deux sommets portant chacun une boucle sur lui-même.

  \item Écrire deux fonctions de comptage des flèches, une par représentation,
    et vérifier qu'elles coïncident sur les quatre graphes.

  \item Répondre par écrit, en commentaire : sur le graphe complet à quatre
    sommets, combien de cases la matrice occupe-t-elle, et combien de valeurs
    les listes stockent-elles ? Même question sur le graphe sans flèche.
    Conclure en une phrase.
\end{enumerate}

\subsection*{Ce qui est évalué}

\begin{center}
\begin{tabular}{|l|c|}
\hline
\textbf{Critère} & \textbf{Points} \\
\hline
Le jeu de tests fourni passe entièrement & $7$ \\
\hline
Les trois préconditions sont refusées & $4$ \\
\hline
L'aller-retour est vérifié sur les quatre graphes & $4$ \\
\hline
Les deux comptages coïncident & $2$ \\
\hline
La réponse écrite donne les quatre nombres et conclut & $3$ \\
\hline
\end{tabular}
\end{center}

\subsection*{Réponse attendue à la question 5}

Graphe complet à quatre sommets : la matrice occupe $16$ cases, les listes
stockent $12$ valeurs. Graphe sans flèche : la matrice occupe toujours $16$
cases, les listes n'en stockent aucune.

La matrice occupe $n^2$ cases quel que soit le nombre de flèches ; les listes en
occupent autant qu'il y a de flèches. L'écart est d'autant plus grand que le
graphe est creux.

\subsection*{Critique}

\begin{itemize}
  \item \textbf{Omettre les sommets isolés.} Une compréhension qui ne garde que
    les sommets ayant au moins un successeur casse l'aller-retour : la matrice
    reconstruite est trop petite. Le graphe sans flèche est le test qui le
    révèle.
  \item \textbf{Les boucles sur soi-même.} La diagonale peut valoir $1$ ; rien
    ne l'interdit, et le quatrième graphe est là pour qu'un candidat qui la
    filtre s'en aperçoive.
  \item \textbf{Vérifier l'aller-retour dans un seul sens.} Reconstruire la
    matrice à partir des listes ne suffit pas : il faut comparer à la matrice
    de départ, pas à une matrice quelconque.
  \item \textbf{Compter les entrées du dictionnaire.} Elles valent toujours $n$.
    Ce qui mesure la mémoire est le nombre total de successeurs stockés.
\end{itemize}

\subsection*{Portée pédagogique}

La situation tient en deux fonctions courtes, et tout l'intérêt est dans les
cas de test : le graphe vide, le graphe complet, les boucles. C'est en les
traversant qu'un candidat découvre les hypothèses implicites de son code.

L'aller-retour fournit une vérification que le candidat peut mener seul, sans
corrigé : si la matrice reconstruite diffère, l'une des deux fonctions est
fausse. La question 5 relie enfin le choix de représentation à une quantité
mesurée, pas à une préférence.
